\begin{lemma}[Weighted subset count by marked intersection]
Let $n,s$ be natural numbers with $s\le n$, let $M\subseteq\{0,\ldots,n-1\}$,
and let $r$ be real. Write $t=|M|$ and assign weight $r$ to each member of
$M$ and weight $1$ to each remaining index. Then
\[
 \sum_{\substack{A\subseteq\{0,\ldots,n-1\}\\|A|=s}}
   \prod_{i\in A}\begin{cases}r&i\in M,\\1&i\notin M\end{cases}
 =\sum_{x=0}^n
   \begin{cases}
    \binom{t}{x}\binom{n-t}{s-x}r^x,&x\le s,\\
    0,&x>s.
   \end{cases}
\]
The empty product and $r^0$ are both $1$.
\end{lemma}
\begin{proof}
Put $U=\{0,\ldots,n-1\}$. For each $s$-subset $A$ of $U$, its weight is
$r^{|A\cap M|}$, since precisely the factors indexed by $A\cap M$ are
$r$, and the other factors are $1$. Its intersection size lies between
$0$ and $s\le n$, so partitioning all $s$-subsets according to this size
accounts for each summand exactly once.

Fix $x\le s$. The map
\[
 A\longmapsto (A\cap M,A\setminus M)
\]
is a bijection from the $s$-subsets having marked intersection size $x$
to the pairs $(B,C)$ where $B\subseteq M$, $|B|=x$,
$C\subseteq U\setminus M$, and $|C|=s-x$.
Indeed, for an $A$ in this fiber, $A\cap M$ has size $x$, and the disjoint
decomposition $A=(A\cap M)\cup(A\setminus M)$ gives
$|A\setminus M|=s-x$. Conversely, for such a pair, $B$ and $C$ are
disjoint, so $A=B\cup C$ has size $x+(s-x)=s$; moreover
$A\cap M=B$ and $A\setminus M=C$. Union is therefore a two-sided inverse.
There are $\binom{t}{x}$ choices for $B$ and
$\binom{n-t}{s-x}$ choices for $C$, by the finite-subset counting formula.
All members of the fiber have weight $r^x$, giving the stated contribution.
If either requested subset size exceeds the corresponding ambient size,
its family is empty and the same counting formula gives zero.

For $x>s$ the fiber is empty because $A\cap M\subseteq A$ implies
$|A\cap M|\le |A|=s$. Its contribution is thus zero, independently of
any convention for truncated subtraction. Summing the fiber contributions
over $x=0,\ldots,n$ proves the identity, also when $s=0$ or $r=0$.
\end{proof}
